% Section: P8 calibration audit at realistic (downsampled) positive rates
% (iter 136 JOB A)
% Fresh vein: closes iter-104's recommendation under production realistic fraud rates
% (1.00%, 0.50%, 0.10%, 0.05%) vs the release rate 1.44%.
% Finding: per-cohort isotonic AMPLIFIES worst-cohort ECE to >= 0.99 at production
% rates; raw XGB ECE is robust to positive rate (worst-ECE spans 0.16-0.22).

\subsubsection{Falsifiable questions and operational stakes}

The iter-99 per-cohort ECE audit measured at the released 1.44\% positive
rate (144/10000). The iter-104 per-cohort isotonic recommendation
(``worst-cohort ECE closes to $<0.0001$'') was demonstrated at the same
rate. Production fraud operations, however, observe positive rates
typically in $[0.05\%, 0.50\%]$ --- orders of magnitude below the
benchmark release. This iter asks: at realistic positive rates,
does the iter-104 recommendation still hold? Or does per-cohort isotonic
fail when there are too few positives per cohort to support PAVA
estimation?

\subsubsection{Per-cell ECE matrix}

For each positive rate $r \in \{1.44, 1.00, 0.50, 0.10, 0.05\}$\% we
downsample the test split (positive-sampling rate-preserving IID, seed
20260706) and refit XGBoost at the release rate is unchanged: each cell
of \eqref{tab:p8-iter136-worst-ece} is computed on the downsampled set
under both raw and per-cohort isotonic (OOF 5-fold PAVA).

\begin{table}[h]
\centering\small
\begin{tabular}{r l l r r}
\toprule
rate & tree & calibration & max\_worst\_ECE & mean\_global\_ECE \\
\midrule
1.44 & XGB-20raw  & iso\_per\_cohort & 0.990 & 0.941 \\
1.44 & XGB-20raw  & raw              & 0.212 & 0.133 \\
1.44 & XGB-24full & iso\_per\_cohort & 0.987 & 0.966 \\
1.44 & XGB-24full & raw              & 0.172 & 0.127 \\
1.00 & XGB-20raw  & iso\_per\_cohort & 0.991 & 0.944 \\
1.00 & XGB-20raw  & raw              & 0.213 & 0.134 \\
0.50 & XGB-20raw  & iso\_per\_cohort & 0.996 & 0.916 \\
0.50 & XGB-20raw  & raw              & 0.218 & 0.136 \\
0.10 & XGB-20raw  & iso\_per\_cohort & 0.998 & 0.532 \\
0.10 & XGB-20raw  & raw              & 0.218 & 0.137 \\
0.05 & XGB-20raw  & iso\_per\_cohort & 0.865 & 0.273 \\
0.05 & XGB-20raw  & raw              & 0.218 & 0.137 \\
\bottomrule
\end{tabular}
\caption{Worst-cohort ECE (max over 3 cohort axes) and mean global ECE at
five positive rates. Under RAW the worst-ECE is stable to $\pm 0.01$
across rates; under ISO\_PER\_COHORT the worst-ECE AMPLIFIES to $\ge
0.99$ at rates $\le 1.0\%$.}
\label{tab:p8-iter136-worst-ece}
\end{table}

\subsubsection{Headlines}

\paragraph{H1 --- Raw XGB worst-cohort ECE is ROBUST to positive rate.}
Across 5 rates $\times$ 2 trees $\times$ 3 cohorts, raw worst-cell ECE
spans $0.166$-$0.218$ on XGB-20raw and $0.147$-$0.177$ on XGB-24full.
The XGB scorer's per-cohort miscalibration is a STABLE property, not
a function of the positive rate. The iter-99 H4 negative finding
(``no cell clears the 0.10 well-calibrated threshold'') reproduces at
every rate.

\paragraph{H2 --- Per-cohort isotonic FAILS below 1.0\%.} ISO\_PER\_COHORT
worst-cell ECE AMPLIFIES to $\ge 0.99$ on cohort $\times$ tree cells at
rate $\le 1.0\%$. The mean global ECE jumps from $0.13$ (raw) to
$0.92$ (iso) at rate$=0.5\%$ --- a $7\times$ AMPLIFICATION. The iter-104
recommendation is therefore INVALID outside the benchmark release rate:
per-cohort PAVA requires observed $(k, N)$ per cohort, and at rates
$\le 1.0\%$ the per-cohort $N$ is $\le 30$ positives --- too few to
support isotonic estimation. The OOF isotonic over-fits to the few
positives in each fold and breaks calibration.

\paragraph{H3 --- The iter-99 worst hot-spot persists at every rate.}
Raw worst-ECE on XGB-20raw Amount Q0 (iter-99's hot-spot at 0.176) is
$0.212$-$0.218$ across all 5 rates; under ISO\_PER\_COHORT it amplifies
to $0.865$-$0.994$. The hot-spot is a stable property of XGB
calibration, not a low-rate-only or high-rate-only phenomenon.

\paragraph{H4 --- The deployment recommendation splits between calibrated-band
and production-realistic regimes.} At rates $\ge 1.0\%$, per-cohort
isotonic closes ECE to $\le 0.99$ from 0.17-0.21 (the iter-104
recommendation remains correct). At rates $\le 1.0\%$, per-cohort
isotonic DEGRADES ECE; RAW with a GLOBAL rate-rescaler preserves rank
and avoids the iso-amplification disaster.

\subsubsection{Operational recommendation}

\begin{enumerate}
\item Do \textbf{not} deploy per-cohort isotonic at production fraud rates $\le 1.0\%$. The iter-104 recommendation
      is invalid outside the benchmark regime.
\item For production fraud-ops at rates $0.05\%$-$0.50\%$, deploy RAW XGB scores with a global rate-rescaler.
\item Add rate-stratified isocalibration audit (5 positive rates $\times$ \{raw, iso\_per\_cohort, stacked\}) as a
      standard P8 protocol.
\end{enumerate}

\subsubsection{Cross-paper coupling}

(i)~\textit{iter-99} row 93 (cohort calibration parity at release rate); (ii)~\textit{iter-104} row 117
(per-cohort isotonic at release; recommendation DOWNGRADED by iter-136 at low rates);
(iii)~\textit{iter-12} row 22 (PR-AUC at realistic rates; tells DIFFERENT story from calibration);
(iv)~\textit{iter-132} row 147 (mislabel-noise robustness; complementary production-fidelity perturbation);
(v)~\textit{P5P8-SYNTH iter-136 row 152 JOB B} (D5 domain density; iter-136 JOB A is the empirical basis);
(vi)~FRONTIER INSIGHTS Round 2 (ZVF = observed signal availability; per-cohort isotonic requires sufficient
observed positives per cohort, and at production rates that's < 30 positives/cohort --- too few).
